\section{Ricerca predittiva}
\label{sec:predictive}
Il metodo B usa il modello per scegliere quali continuazioni esplorare prima.
Per confrontarle deve considerare sia la distanza già percorsa, $g(s)$,
sia quella ancora prevista, $\widehat h(s)$. Chiamo $\widehat f(s)$ la loro
somma, cioè la stima della distanza totale del giro che passa da quello stato:
\[
 \widehat f(s)=g(s)+\widehat h(s).
\]
Ad esempio, uno stato con $g=100$ e $\widehat h=80$ ha costo totale previsto
$180$. Uno con $g=140$ e $\widehat h=50$ ha costo totale previsto $190$.
Il primo viene preferito, anche se il secondo ha una previsione residua minore.

Dopo ogni espansione, B ordina i nuovi candidati secondo questa stima,
in modo da proseguire prima lungo il meno costoso. Il confronto avviene
fra i successori dello stesso stato: non viene riordinata tutta la frontiera.
La ricerca conserva quindi il comportamento in profondità descritto per A.

Questa precisazione distingue B da A*: pur usando anch'esso una somma di
costo percorso e costo previsto, A* sceglie il prossimo stato confrontando
globalmente i candidati della frontiera. Nel progetto l'ordinamento è locale.

A e B usano la stessa regola di potatura; cambia l'ordine di esplorazione.
Il confronto A--B misura dunque il contributo della previsione. C aggiunge
la regola di scelta descritta nella sezione seguente, quindi B--C permette
di misurare il contributo aggiuntivo del regolatore.

Una previsione sbagliata può far esplorare tardi un buon ramo e peggiorare
il giro trovato entro il tempo disponibile. Non rende però scorretta la
potatura, che continua a usare soltanto la distanza già percorsa.
Se la ricerca riesce a esaminare o escludere correttamente tutte le
alternative, trova comunque l'ottimo. I pesi del modello restano fissi
per tutta la ricerca.
